\documentclass[a4paper,11pt]{article}
\usepackage[utf8]{inputenc}
\usepackage[french]{babel}
\usepackage[margin=2.5cm]{geometry}
\usepackage{hyperref}
\usepackage{xcolor}
\usepackage{tcolorbox}
\usepackage{enumitem}
\usepackage{fontawesome5}
\usepackage{graphicx}

\definecolor{vscodeblue}{RGB}{0,122,204}
\definecolor{gitgreen}{RGB}{40,167,69}
\definecolor{commandbg}{RGB}{245,245,245}

\title{\textbf{Guide de Synchronisation Git avec VS Code}\\
\large Projet immunomath}
\author{Équipe IMMUN4CURE}
\date{Octobre 2025}

\begin{document}

\maketitle

\section*{Introduction}

Ce guide explique comment synchroniser votre branche locale \texttt{main} avec le dépôt GitHub du projet \textbf{immunomath} en utilisant Visual Studio Code.

\vspace{0.5cm}

\begin{tcolorbox}[colback=vscodeblue!5,colframe=vscodeblue,title={\faInfoCircle\ Objectif}]
Récupérer les dernières modifications du projet depuis GitHub pour travailler sur la version la plus récente du code.
\end{tcolorbox}

\section{Méthode 1 : Interface Graphique VS Code \faMousePointer}

\subsection*{Étape 1 : Vérifier/Changer de branche}

\begin{enumerate}[leftmargin=2cm]
    \item Regardez en \textbf{bas à gauche} de la fenêtre VS Code
    \item Vous verrez le nom de la branche actuelle (ex: \texttt{main}, \texttt{feature/xyz})
    \item Cliquez sur ce nom de branche
    \item Dans le menu déroulant qui apparaît, sélectionnez \texttt{\textbf{main}}
\end{enumerate}

\begin{tcolorbox}[colback=commandbg,colframe=gray,leftrule=3mm]
\faLightbulb\ \textbf{Astuce :} Cette action équivaut à la commande \texttt{git checkout main}
\end{tcolorbox}

\subsection*{Étape 2 : Synchroniser avec GitHub}

Vous avez \textbf{deux options} :

\subsubsection*{Option A : Bouton de synchronisation (recommandé)}

\begin{enumerate}[leftmargin=2cm]
    \item Cherchez l'icône de \textbf{synchronisation} \faSync\ dans la barre d'état (en bas de VS Code)
    \item OU cherchez les flèches circulaires en haut à droite de l'onglet Source Control
    \item Cliquez sur cette icône
    \item VS Code va automatiquement :
    \begin{itemize}
        \item[\faDownload] Télécharger les modifications distantes
        \item[\faMagic] Fusionner les changements avec votre version locale
        \item[\faUpload] Pousser vos modifications locales (si nécessaire)
    \end{itemize}
\end{enumerate}

\subsubsection*{Option B : Menu Palette de commandes}

\begin{enumerate}[leftmargin=2cm]
    \item Appuyez sur \texttt{\textbf{Ctrl+Shift+P}} (Windows/Linux) ou \texttt{\textbf{Cmd+Shift+P}} (Mac)
    \item Tapez : \texttt{Git: Pull}
    \item Appuyez sur \texttt{Entrée}
\end{enumerate}

\vspace{0.5cm}

\begin{tcolorbox}[colback=gitgreen!10,colframe=gitgreen,title={\faCheckCircle\ Confirmation}]
Un message apparaîtra en bas à droite confirmant la synchronisation réussie.
\end{tcolorbox}

\newpage

\section{Méthode 2 : Terminal Intégré \faTerminal}

Si vous préférez utiliser les commandes Git directement :

\subsection*{Étape 1 : Ouvrir le terminal}

\begin{enumerate}[leftmargin=2cm]
    \item Appuyez sur \texttt{\textbf{Ctrl+`}} (touche accent grave sous Échap)
    \item OU Menu \fabars\ $\rightarrow$ Terminal $\rightarrow$ Nouveau terminal
\end{enumerate}

\subsection*{Étape 2 : Exécuter les commandes}

Tapez les commandes suivantes dans le terminal :

\begin{tcolorbox}[colback=black,colframe=black,colupper=white,fontupper=\ttfamily]
\$ git checkout main\\
\$ git pull origin main
\end{tcolorbox}

\vspace{0.3cm}

\textbf{Explication :}
\begin{itemize}
    \item \texttt{git checkout main} : Bascule sur la branche principale
    \item \texttt{git pull origin main} : Récupère et fusionne les modifications de GitHub
\end{itemize}

\section{Gestion des Conflits \faExclamationTriangle}

Si des conflits apparaissent (vous et un collaborateur avez modifié les mêmes lignes) :

\begin{enumerate}[leftmargin=2cm]
    \item VS Code ouvrira automatiquement l'éditeur de conflits
    \item Vous verrez des sections marquées :
    \begin{itemize}
        \item \textcolor{blue}{\texttt{<<<<<<< HEAD}} : Vos modifications locales
        \item \textcolor{red}{\texttt{=======}} : Séparateur
        \item \textcolor{green}{\texttt{>>>>>>> origin/main}} : Modifications distantes
    \end{itemize}
    \item Utilisez les boutons au-dessus du conflit :
    \begin{itemize}
        \item \texttt{Accept Current Change} : Garder vos modifications
        \item \texttt{Accept Incoming Change} : Garder les modifications de GitHub
        \item \texttt{Accept Both Changes} : Garder les deux
        \item \texttt{Compare Changes} : Voir les différences côte à côte
    \end{itemize}
    \item Sauvegardez le fichier après résolution
    \item Validez la fusion via l'interface Source Control
\end{enumerate}

\section{Bonnes Pratiques \faCheckDouble}

\begin{enumerate}
    \item \textbf{Synchronisez régulièrement} : Au début de chaque session de travail
    \item \textbf{Commitez vos changements} avant de synchroniser pour éviter les pertes
    \item \textbf{Vérifiez la branche active} avant de commencer à travailler
    \item \textbf{Communiquez} avec l'équipe sur les fichiers modifiés pour éviter les conflits
\end{enumerate}

\section{Aide Rapide \faQuestionCircle}

\begin{center}
\begin{tabular}{|p{6cm}|p{8cm}|}
\hline
\textbf{Problème} & \textbf{Solution} \\
\hline
Je ne vois pas l'icône de synchronisation & Vérifiez que vous êtes dans un dépôt Git. Ouvrez le panneau Source Control (Ctrl+Shift+G) \\
\hline
Message d'erreur lors du pull & Assurez-vous d'avoir committé ou mis de côté vos modifications locales \\
\hline
Conflit de fusion & Suivez les instructions de la Section 3 pour résoudre les conflits \\
\hline
Branche main introuvable & Vérifiez le nom exact de la branche (peut être \texttt{master} sur certains dépôts) \\
\hline
\end{tabular}
\end{center}

\vspace{1cm}

\begin{tcolorbox}[colback=yellow!10,colframe=orange,title={\faPhone\ Contact}]
Pour toute question ou problème, n'hésitez pas à contacter l'équipe technique du projet IMMUN4CURE.
\end{tcolorbox}

\vspace{0.5cm}

\begin{center}
\textit{Dépôt GitHub : } \url{https://github.com/azerad/immunomath}
\end{center}

\end{document}
